\documentclass[11pt,a4paper]{article}

% -------------------------------------------------
% Packages
% -------------------------------------------------
\usepackage[margin=2.2cm]{geometry}
\usepackage{kotex}
\usepackage{amsmath,amssymb}
\usepackage{booktabs}
\usepackage{array}
\usepackage{xcolor}
\usepackage{listings}
\usepackage[most]{tcolorbox}
\usepackage{enumitem}
\usepackage{hyperref}
\usepackage{fancyhdr}
\usepackage{titlesec}

% -------------------------------------------------
% Colors
% -------------------------------------------------
\definecolor{codebg}{RGB}{247,247,247}
\definecolor{codeframe}{RGB}{210,210,210}
\definecolor{keywordcolor}{RGB}{0,70,160}
\definecolor{commentcolor}{RGB}{0,120,80}
\definecolor{stringcolor}{RGB}{170,50,50}
\definecolor{titleblue}{RGB}{35,75,130}

% -------------------------------------------------
% Hyperlinks
% -------------------------------------------------
\hypersetup{
    colorlinks=true,
    linkcolor=titleblue,
    urlcolor=titleblue
}

% -------------------------------------------------
% Header / Footer
% -------------------------------------------------
\pagestyle{fancy}
\fancyhf{}
\lhead{Python Programming}
\rhead{Day 8}
\cfoot{\thepage}
\setlength{\headheight}{14pt}

% -------------------------------------------------
% Section style
% -------------------------------------------------
\titleformat{\section}
  {\Large\bfseries\color{titleblue}}
  {\thesection.}{0.6em}{}

\titleformat{\subsection}
  {\large\bfseries}
  {\thesubsection.}{0.5em}{}

% -------------------------------------------------
% Python code style
% -------------------------------------------------
\lstdefinestyle{pythonstyle}{
    language=Python,
    backgroundcolor=\color{codebg},
    frame=single,
    rulecolor=\color{codeframe},
    basicstyle=\ttfamily\small,
    keywordstyle=\color{keywordcolor}\bfseries,
    commentstyle=\color{commentcolor},
    stringstyle=\color{stringcolor},
    numbers=left,
    numberstyle=\tiny\color{gray},
    numbersep=8pt,
    tabsize=4,
    showstringspaces=false,
    breaklines=true,
    columns=fullflexible,
    keepspaces=true
}

\lstset{style=pythonstyle}

% -------------------------------------------------
% Boxes
% -------------------------------------------------
\newtcolorbox{learningbox}{
    colback=blue!4,
    colframe=titleblue,
    title=\textbf{학습 목표},
    fonttitle=\bfseries
}

\newtcolorbox{importantbox}{
    colback=yellow!8,
    colframe=orange!70!black,
    title=\textbf{핵심 포인트},
    fonttitle=\bfseries
}

\newtcolorbox{practicebox}{
    colback=red!3,
    colframe=red!55!black,
    title=\textbf{실습},
    fonttitle=\bfseries
}

% -------------------------------------------------
% Document
% -------------------------------------------------
\begin{document}

\begin{titlepage}
    \centering
    \vspace*{3cm}

    {\Huge\bfseries Python Programming\par}
    \vspace{0.8cm}
    {\LARGE Day 8: Pandas 기초\par}

    \vspace{2cm}

    {\large Series, DataFrame, Indexing, Modification, CSV, Statistics\par}

    \vfill
\end{titlepage}

\tableofcontents
\newpage

% =================================================
\section{학습 목표}
% =================================================

\begin{learningbox}
이 장을 마치면 다음 내용을 설명하고 직접 코드로 작성할 수 있다.
\begin{itemize}[leftmargin=1.5em]
    \item Pandas의 목적과 기본 자료구조 설명하기
    \item \texttt{Series}를 생성하고 인덱스로 값 선택하기
    \item 딕셔너리로 \texttt{DataFrame} 생성하기
    \item 열 선택과 \texttt{iloc}을 이용한 위치 기반 데이터 선택하기
    \item 새로운 열을 추가하고 \texttt{drop()}으로 열 삭제하기
    \item \texttt{to\_csv()}와 \texttt{read\_csv()}로 CSV 저장 및 불러오기
    \item \texttt{mean()}, \texttt{sum()}, \texttt{max()}, \texttt{min()}, \texttt{describe()} 사용하기
\end{itemize}
\end{learningbox}

% =================================================
\section{Pandas와 Series}
% =================================================

Pandas는 Python에서 표 형태의 데이터를 저장하고 분석하기 위한 라이브러리이다. NumPy가 수치 배열 계산에 강하다면, Pandas는 CSV와 같은 실제 데이터 파일을 읽고 원하는 행과 열을 선택하며 통계를 계산하는 작업에 적합하다.

Pandas를 설치하지 않았다면 터미널에서 다음 명령을 한 번 실행한다.

\begin{verbatim}
pip install pandas
\end{verbatim}

Pandas는 일반적으로 \texttt{pd}라는 별칭으로 불러온다.

\begin{lstlisting}
import pandas as pd
\end{lstlisting}

\subsection{Series}

\texttt{Series}는 값과 인덱스를 함께 가지는 1차원 자료구조이다. 인덱스를 따로 지정하지 않으면 \texttt{0, 1, 2, ...}가 자동으로 생성되지만, 문자열이나 다른 값을 직접 인덱스로 지정할 수도 있다.

\begin{lstlisting}
scores = pd.Series(
    [90, 85, 100],
    index=["Math", "English", "Physics"]
)
\end{lstlisting}

특정 인덱스의 값은 대괄호 안에 인덱스를 넣어 선택한다.

\begin{lstlisting}
print(scores["English"])
\end{lstlisting}

\subsection{최종 코드: \texttt{01\_series.py}}

\begin{lstlisting}
import pandas as pd

scores = pd.Series(
    [90, 85, 100],
    index=["Math", "English", "Physics"]
)

print(scores)
print()

print(scores["English"])
print()

print(scores[["Math", "Physics"]])
print()

print(scores.values)
print(scores.index)
\end{lstlisting}

실행 결과는 다음과 같다.

\begin{verbatim}
Math        90
English     85
Physics    100
dtype: int64

85

Math        90
Physics    100
dtype: int64

[ 90  85 100]
Index(['Math', 'English', 'Physics'], dtype='object')
\end{verbatim}

\begin{importantbox}
\texttt{Series}의 인덱스는 \texttt{0, 1, 2, ...}로 고정되지 않는다. 직접 지정한 인덱스를 이용하면 딕셔너리처럼 의미 있는 이름으로 값에 접근할 수 있다.
\end{importantbox}

% =================================================
\section{DataFrame}
% =================================================

\texttt{DataFrame}은 행과 열을 가지는 2차원 자료구조이다. 엑셀 표와 비슷하며, 여러 개의 \texttt{Series}가 열 방향으로 모인 형태로 이해할 수 있다.

DataFrame은 보통 딕셔너리를 이용하여 만든다. 딕셔너리의 key는 열 이름이 되고, value에 들어 있는 리스트가 해당 열의 데이터가 된다.

\begin{center}
\begin{tabular}{lrr}
\toprule
Name & Math & English \\
\midrule
Kim & 90 & 88 \\
Lee & 85 & 91 \\
Park & 100 & 95 \\
\bottomrule
\end{tabular}
\end{center}

각 열의 길이는 같아야 한다. 길이가 다르면 DataFrame을 만들 수 없다.

\subsection{최종 코드: \texttt{02\_dataframe.py}}

\begin{lstlisting}
import pandas as pd

data = {
    "Name": ["Kim", "Lee", "Park"],
    "Math": [90, 85, 100],
    "English": [88, 91, 95]
}

df = pd.DataFrame(data)

print(df)
print()
print(type(df))
\end{lstlisting}

실행 결과는 다음과 같다.

\begin{verbatim}
   Name  Math  English
0   Kim    90       88
1   Lee    85       91
2  Park   100       95

<class 'pandas.core.frame.DataFrame'>
\end{verbatim}

\begin{importantbox}
DataFrame은 열마다 서로 다른 자료형을 가질 수 있다. 예를 들어 이름은 문자열, 점수는 정수, 합격 여부는 불리언으로 저장할 수 있다.
\end{importantbox}

% =================================================
\section{데이터 선택}
% =================================================

DataFrame을 만든 뒤에는 필요한 열, 행 또는 하나의 값을 선택할 수 있다.

\subsection{열 선택}

하나의 열은 열 이름을 대괄호 안에 넣어 선택한다.

\begin{lstlisting}
df["Math"]
\end{lstlisting}

하나의 열을 선택하면 결과는 \texttt{Series}이다. 여러 열을 선택할 때는 열 이름 목록을 전달한다.

\begin{lstlisting}
df[["Math", "English"]]
\end{lstlisting}

여러 열을 선택하면 결과는 \texttt{DataFrame}이다.

\subsection{\texttt{iloc}을 이용한 위치 선택}

\texttt{iloc}은 정수 위치를 기준으로 데이터를 선택한다. 위치는 Python 인덱스와 마찬가지로 0부터 시작한다.

\begin{center}
\begin{tabular}{ll}
\toprule
표현 & 의미 \\
\midrule
\texttt{df.iloc[0]} & 첫 번째 행 \\
\texttt{df.iloc[1]} & 두 번째 행 \\
\texttt{df.iloc[1, 2]} & 두 번째 행, 세 번째 열의 값 \\
\bottomrule
\end{tabular}
\end{center}

\texttt{df.iloc[1, 2]}를 ``1행 2열''이라고 표현하면 일반적인 행렬 표기와 혼동될 수 있다. 따라서 위치를 설명할 때는 ``행 인덱스 1, 열 인덱스 2'' 또는 ``두 번째 행, 세 번째 열''이라고 표현하는 것이 정확하다.

\subsection{최종 코드: \texttt{03\_indexing.py}}

\begin{lstlisting}
import pandas as pd

data = {
    "Name": ["Kim", "Lee", "Park"],
    "Math": [90, 85, 100],
    "English": [88, 91, 95]
}

df = pd.DataFrame(data)

print("=== Math Column ===")
print(df["Math"])
print()

print("=== Math and English ===")
print(df[["Math", "English"]])
print()

print("=== First Row ===")
print(df.iloc[0])
print()

print("=== Second Row, Third Column ===")
print(df.iloc[1, 2])
\end{lstlisting}

실행 결과는 다음과 같다.

\begin{verbatim}
=== Math Column ===
0     90
1     85
2    100
Name: Math, dtype: int64

=== Math and English ===
   Math  English
0    90       88
1    85       91
2   100       95

=== First Row ===
Name       Kim
Math        90
English     88
Name: 0, dtype: object

=== Second Row, Third Column ===
91
\end{verbatim}

\begin{importantbox}
열 이름으로 선택할 때는 \texttt{df["Column"]}을 사용하고, 실제 저장 위치를 기준으로 행과 열을 선택할 때는 \texttt{iloc}을 사용한다.
\end{importantbox}

% =================================================
\section{데이터 추가와 삭제}
% =================================================

새로운 열을 추가할 때는 딕셔너리에 항목을 추가하듯이 열 이름에 값을 대입한다.

\begin{lstlisting}
df["Physics"] = [95, 80, 90]
\end{lstlisting}

같은 이름의 열이 이미 존재하면 기존 값이 새로운 값으로 바뀐다.

열을 삭제할 때는 \texttt{drop()}을 사용한다.

\begin{lstlisting}
df = df.drop(columns=["Math"])
\end{lstlisting}

\texttt{drop()}은 수정된 새로운 DataFrame을 반환하므로, 결과를 다시 \texttt{df}에 저장한다.

\subsection{최종 코드: \texttt{04\_modify.py}}

\begin{lstlisting}
import pandas as pd

data = {
    "Name": ["Kim", "Lee", "Park"],
    "Math": [90, 85, 100]
}

df = pd.DataFrame(data)

print("=== Original ===")
print(df)
print()

df["Physics"] = [95, 80, 90]

print("=== After Adding Physics ===")
print(df)
print()

df = df.drop(columns=["Math"])

print("=== After Deleting Math ===")
print(df)
\end{lstlisting}

실행 결과는 다음과 같다.

\begin{verbatim}
=== Original ===
   Name  Math
0   Kim    90
1   Lee    85
2  Park   100

=== After Adding Physics ===
   Name  Math  Physics
0   Kim    90       95
1   Lee    85       80
2  Park   100       90

=== After Deleting Math ===
   Name  Physics
0   Kim       95
1   Lee       80
2  Park       90
\end{verbatim}

행도 \texttt{drop(index=...)}으로 삭제할 수 있다.

\begin{lstlisting}
df = df.drop(index=1)
\end{lstlisting}

\begin{importantbox}
새로운 열의 데이터 개수는 DataFrame의 행 개수와 같아야 한다. 개수가 다르면 길이가 맞지 않는다는 오류가 발생한다.
\end{importantbox}

% =================================================
\section{CSV 저장과 불러오기}
% =================================================

CSV는 Comma-Separated Values의 약자로, 쉼표를 이용해 표 데이터를 저장하는 텍스트 파일 형식이다. Pandas를 사용하면 DataFrame을 CSV로 저장하고 다시 읽는 작업을 한 줄씩으로 처리할 수 있다.

\subsection{CSV 저장}

\begin{lstlisting}
df.to_csv("students.csv", index=False)
\end{lstlisting}

\texttt{index=False}를 사용하면 DataFrame의 자동 인덱스가 파일에 별도 열로 저장되지 않는다.

\subsection{CSV 불러오기}

\begin{lstlisting}
loaded = pd.read_csv("students.csv")
\end{lstlisting}

\texttt{read\_csv()}의 반환값은 DataFrame이다.

\subsection{최종 코드: \texttt{05\_csv.py}}

\begin{lstlisting}
import pandas as pd

data = {
    "Name": ["Kim", "Lee", "Park"],
    "Math": [90, 85, 100],
    "English": [88, 91, 95]
}

df = pd.DataFrame(data)

df.to_csv("students.csv", index=False)

loaded = pd.read_csv("students.csv")

print("=== Loaded CSV ===")
print(loaded)
print()
print(type(loaded))
\end{lstlisting}

실행 결과는 다음과 같다.

\begin{verbatim}
=== Loaded CSV ===
   Name  Math  English
0   Kim    90       88
1   Lee    85       91
2  Park   100       95

<class 'pandas.core.frame.DataFrame'>
\end{verbatim}

\begin{importantbox}
CSV는 자료형이나 서식을 완전히 보존하는 엑셀 파일과 다르다. 하지만 구조가 단순하고 여러 프로그램에서 쉽게 읽을 수 있어 데이터 분석에서 자주 사용된다.
\end{importantbox}

% =================================================
\section{기초 통계}
% =================================================

Pandas의 Series와 DataFrame은 기본적인 통계 함수를 제공한다.

\begin{center}
\begin{tabular}{ll}
\toprule
함수 & 기능 \\
\midrule
\texttt{mean()} & 평균 \\
\texttt{sum()} & 합계 \\
\texttt{max()} & 최댓값 \\
\texttt{min()} & 최솟값 \\
\texttt{describe()} & 주요 기술통계 요약 \\
\bottomrule
\end{tabular}
\end{center}

\texttt{describe()}는 데이터 개수, 평균, 표준편차, 최솟값, 사분위수, 최댓값을 한 번에 계산한다.

\begin{center}
\begin{tabular}{ll}
\toprule
항목 & 의미 \\
\midrule
\texttt{count} & 데이터 개수 \\
\texttt{mean} & 평균 \\
\texttt{std} & 표본 표준편차 \\
\texttt{min} & 최솟값 \\
\texttt{25\%} & 제1사분위수 \\
\texttt{50\%} & 중앙값 \\
\texttt{75\%} & 제3사분위수 \\
\texttt{max} & 최댓값 \\
\bottomrule
\end{tabular}
\end{center}

\subsection{최종 코드: \texttt{06\_statistics.py}}

\begin{lstlisting}
import pandas as pd

data = {
    "Math": [90, 85, 100, 70, 95]
}

df = pd.DataFrame(data)

print("Mean:", df["Math"].mean())
print("Sum:", df["Math"].sum())
print("Max:", df["Math"].max())
print("Min:", df["Math"].min())

print("\n=== Describe ===")
print(df.describe())
\end{lstlisting}

실행 결과는 다음과 같다.

\begin{verbatim}
Mean: 88.0
Sum: 440
Max: 100
Min: 70

=== Describe ===
             Math
count    5.000000
mean    88.000000
std     11.510864
min     70.000000
25%     85.000000
50%     90.000000
75%     95.000000
max    100.000000
\end{verbatim}

\begin{importantbox}
통계 함수는 숫자 열을 선택한 뒤 적용하는 것이 기본이다. CSV에서 읽어온 DataFrame에도 같은 방식으로 사용할 수 있다.
\end{importantbox}

% =================================================
\section{Day 8 종합 실습: 학생 성적 관리}
% =================================================

이번 종합 실습에서는 DataFrame 생성, 데이터 선택, 열 추가 및 삭제, CSV 저장과 불러오기, 기초 통계 기능을 하나의 프로그램에서 연결한다.

\subsection{프로그램 요구사항}

\begin{enumerate}[leftmargin=1.8em]
    \item 네 학생의 이름, 수학 점수, 영어 점수로 DataFrame을 생성한다.
    \item 물리 점수 열을 추가한다.
    \item 전체 DataFrame과 수학 점수 열을 출력한다.
    \item \texttt{iloc}을 이용해 두 번째 학생의 정보를 출력한다.
    \item 영어 점수 열을 삭제한 뒤 DataFrame을 출력한다.
    \item DataFrame을 \texttt{student.csv}에 저장하고 다시 불러온다.
    \item 불러온 데이터의 수학 평균, 최댓값, 최솟값을 출력한다.
    \item \texttt{describe()}로 주요 통계 정보를 출력한다.
\end{enumerate}

\subsection{최종 코드: \texttt{practice/practice\_day08.py}}

\begin{lstlisting}
import pandas as pd

data = {
    "Name": ["Kim", "Lee", "Park", "Choi"],
    "Math": [90, 85, 100, 95],
    "English": [88, 91, 95, 87]
}

df = pd.DataFrame(data)

df["Physics"] = [92, 84, 98, 90]

print("==== Student Data ====")
print(df)

print()
print(df["Math"])

print()
print(df.iloc[1])

print()
df = df.drop(columns=["English"])
print(df)

df.to_csv("student.csv", index=False)
loaded = pd.read_csv("student.csv")

print()
print("Average :", loaded["Math"].mean())
print("Maximum :", loaded["Math"].max())
print("Minimum :", loaded["Math"].min())

print()
print(loaded.describe())
\end{lstlisting}

\subsection{실행 결과}

\begin{verbatim}
==== Student Data ====
   Name  Math  English  Physics
0   Kim    90       88       92
1   Lee    85       91       84
2  Park   100       95       98
3  Choi    95       87       90

0     90
1     85
2    100
3     95
Name: Math, dtype: int64

Name       Lee
Math        85
English     91
Physics     84
Name: 1, dtype: object

   Name  Math  Physics
0   Kim    90       92
1   Lee    85       84
2  Park   100       98
3  Choi    95       90

Average : 92.5
Maximum : 100
Minimum : 85

             Math    Physics
count    4.000000   4.000000
mean    92.500000  91.000000
std      6.454972   5.887841
min     85.000000  84.000000
25%     88.750000  88.500000
50%     92.500000  91.000000
75%     96.250000  93.500000
max    100.000000  98.000000
\end{verbatim}

\subsection{프로그램 구조 분석}

\begin{center}
\begin{tabular}{ll}
\toprule
기능 & 사용한 개념 \\
\midrule
학생 데이터 생성 & \texttt{pd.DataFrame()} \\
물리 점수 추가 & 새 열 대입 \\
수학 점수 선택 & 열 인덱싱 \\
두 번째 학생 선택 & \texttt{iloc} \\
영어 점수 삭제 & \texttt{drop()} \\
파일 저장과 불러오기 & \texttt{to\_csv()}, \texttt{read\_csv()} \\
점수 통계 계산 & \texttt{mean()}, \texttt{max()}, \texttt{min()} \\
전체 통계 확인 & \texttt{describe()} \\
\bottomrule
\end{tabular}
\end{center}

\begin{importantbox}
CSV를 다시 읽은 뒤에는 \texttt{loaded}를 사용하여 통계를 계산했다. 이렇게 하면 저장된 파일이 정상적으로 불러와졌는지까지 함께 확인할 수 있다.
\end{importantbox}

% =================================================
\section{실습 파일 목록}
% =================================================

\begin{practicebox}
Day 8의 학습 파일은 다음과 같이 관리한다.
\begin{itemize}[leftmargin=1.5em]
    \item \texttt{01\_series.py}
    \item \texttt{02\_dataframe.py}
    \item \texttt{03\_indexing.py}
    \item \texttt{04\_modify.py}
    \item \texttt{05\_csv.py}
    \item \texttt{06\_statistics.py}
    \item \texttt{practice/practice\_day08.py}
    \item \texttt{python\_day08.tex}
    \item \texttt{python\_day08.pdf}
\end{itemize}
\end{practicebox}

% =================================================
\section{핵심 요약}
% =================================================

\begin{itemize}[leftmargin=1.5em]
    \item Pandas는 표 형태의 데이터를 저장하고 분석하기 위한 Python 라이브러리이다.
    \item \texttt{Series}는 값과 인덱스를 가지는 1차원 자료구조이다.
    \item Series의 인덱스는 자동 숫자 인덱스뿐 아니라 문자열 등으로 직접 지정할 수 있다.
    \item \texttt{DataFrame}은 행과 열을 가지는 2차원 표 형태의 자료구조이다.
    \item 딕셔너리의 key는 DataFrame의 열 이름이 되고 value는 열 데이터가 된다.
    \item \texttt{df["Column"]}은 하나의 열을 Series로 반환한다.
    \item \texttt{df[["A", "B"]]}는 여러 열을 DataFrame으로 반환한다.
    \item \texttt{iloc}은 0부터 시작하는 정수 위치를 기준으로 데이터를 선택한다.
    \item \texttt{df.iloc[1, 2]}는 두 번째 행, 세 번째 열의 값을 의미한다.
    \item 새 열은 \texttt{df["New"] = values}로 추가하거나 수정한다.
    \item \texttt{drop(columns=[...])}은 지정한 열을 삭제한 새로운 DataFrame을 반환한다.
    \item \texttt{to\_csv()}는 DataFrame을 CSV로 저장한다.
    \item \texttt{read\_csv()}는 CSV를 DataFrame으로 불러온다.
    \item \texttt{index=False}를 사용하면 자동 인덱스가 CSV에 저장되지 않는다.
    \item \texttt{mean()}, \texttt{sum()}, \texttt{max()}, \texttt{min()}으로 기본 통계를 계산한다.
    \item \texttt{describe()}는 숫자 열의 주요 기술통계를 한 번에 보여준다.
    \item Pandas의 고급 기능은 실제 데이터 분석이나 프로젝트에서 필요한 시점에 추가로 학습한다.
\end{itemize}

\end{document}
